\documentclass[a4paper,10pt,sans]{moderncv}
%----------------------------------------------------------------------------------------
%   PACKAGES
%----------------------------------------------------------------------------------------

% Programming facilities
\usepackage{etoolbox}
\usepackage{ifxetex}
\usepackage{ifluatex}

% Encoding
\usepackage[T1]{fontenc}
\ifboolexpr{bool{xetex} or bool{luatex}}{%
    \usepackage{fontspec}
}{%
    \usepackage[utf8]{inputenc}
}

%\usepackage[scale=0.975]{geometry}
\usepackage[top=0.85cm, bottom=0.85cm, left=0.85cm, right=0.85cm]{geometry}
\recomputelengths{}

\usepackage{graphicx}
\usepackage[firstyear=2013,lastyear=2025]{moderntimeline}
\usepackage{datenumber,fp}
\usepackage{fontawesome5}
\usepackage[bibstyle=ieee,citestyle=numeric-comp,language=english]{biblatex}
\usepackage{relsize}
\usepackage{hyperref}
%\usepackage{mnsymbol}
%\usepackage{academicons}

%----------------------------------------------------------------------------------------
%   SETTINGS
%----------------------------------------------------------------------------------------

%\makeatletter % changes the catcode of @ to 11
%modify here packages internal commands
%\makeatother % changes the catcode of @ back to 12

\makeatletter

%-------------------------------------------
%   Add preinfo to header
%-------------------------------------------

\newcommand*{\preinfo}[1]{\def\@preinfo{#1}}

% as \quotewidth can't be changed, we define \quotecustomwidth
\newlength{\quotecustomwidth}
\setlength{\quotecustomwidth}{0.65\textwidth}% default value
% optional maketitle width to force a certain width (if set to 0pt, the width is calculated automatically)
\renewcommand*{\makecvtitle}{%
	{\hypersetup{urlcolor=color2}
	% recompute lengths (in case we are switching from letter to resume, or vice versa)
	\recomputecvlengths%
	% optional detailed information (pre-rendering)
	\def\phonesdetails{}%
	\collectionloop{phones}{% the key holds the phone type (=symbol command prefix), the item holds the number
		\protected@edef\phonesdetails{\phonesdetails\protect\makenewline\csname\collectionloopkey phonesymbol\endcsname\collectionloopitem}}%
	\def\socialsdetails{}%
	\collectionloop{socials}{% the key holds the social type (=symbol command prefix), the item holds the link
		\protected@edef\socialsdetails{\socialsdetails\protect\makenewline\csname\collectionloopkey socialsymbol\endcsname\collectionloopitem}}%
	\newbox{\makecvtitledetailsbox}%
	\savebox{\makecvtitledetailsbox}{%
		\addressfont\color{color2}%
		\begin{tabular}[b]{@{}r@{}}%
			\ifthenelse{\isundefined{\@preinfo}}{}{\makenewline{\@preinfo}}%
			\ifthenelse{\isundefined{\@addressstreet}}{}{\makenewline\addresssymbol\@addressstreet%
			\ifthenelse{\equal{\@addresscity}{}}{}{\makenewline\@addresscity}% if \addresstreet is defined, \addresscity and addresscountry will always be defined but could be empty
			\ifthenelse{\equal{\@addresscountry}{}}{}{\makenewline\@addresscountry}}%
			\phonesdetails% needs to be pre-rendered as loops and tabulars seem to conflict
			\ifthenelse{\isundefined{\@email}}{}{\makenewline\emailsymbol\emaillink{\@email}}%
			\ifthenelse{\isundefined{\@homepage}}{}{\makenewline\homepagesymbol\httplink{\@homepage}}%
			\socialsdetails% needs to be pre-rendered as loops and tabulars seem to conflict
			\ifthenelse{\isundefined{\@extrainfo}}{}{\makenewline\@extrainfo}%
		\end{tabular}
	}%
	% optional photo (pre-rendering)
	\newbox{\makecvtitlepicturebox}%
	\savebox{\makecvtitlepicturebox}{%
		\ifthenelse{\isundefined{\@photo}}%
			{}%
			{%
				\hspace*{\separatorcolumnwidth}%
				\color{color1}%
				\setlength{\fboxrule}{\@photoframewidth}%
				\ifdim\@photoframewidth=0pt%
					\setlength{\fboxsep}{0pt}\fi%
				\framebox{\includegraphics[width=\@photowidth]{\@photo}}%
			}%
	}%
	% name and title
	\newlength{\makecvtitledetailswidth}\settowidth{\makecvtitledetailswidth}{\usebox{\makecvtitledetailsbox}}%
	\newlength{\makecvtitlepicturewidth}\settowidth{\makecvtitlepicturewidth}{\usebox{\makecvtitlepicturebox}}%
	\begin{minipage}[b]{\textwidth-\makecvtitledetailswidth-\makecvtitlepicturewidth - 10mm}%
		\raggedright
		\titlestyle{\@title}\\
		\hfill\\
		\namestyle{\@firstname\ \@lastname}
	\end{minipage}%
	\hfill%
	% optional detailed information (rendering)
	\llap{\usebox{\makecvtitledetailsbox}}% \llap is used to suppress the width of the box, allowing overlap if the value of makecvtitlenamewidth is forced
	% optional photo (rendering)
	\usebox{\makecvtitlepicturebox}\\[2.5em]%
	% optional quote
	\ifthenelse{\isundefined{\@quote}}%
		{}%
		{{\centering\begin{minipage}{\quotecustomwidth}\vspace*{\quoteDeletedSpace}\centering\quotestyle{\@quote}\end{minipage}\\[2.5em]}}%
	\par% to avoid weird spacing bug at the first section if no blank line is left after \makecvtitle
	}
}

%-------------------------------------------
%   Change tllabelcventry label position
%-------------------------------------------

\renewcommand{\tllabelcventry}[9][color1]{%
	\tl@formatendyear{#3}
	\tl@formatstartyear{#2}
	\cventry{
		\tikz[baseline=0pt]{
			\fill [\tl@runningcolor] (0,0) rectangle (\hintscolumnwidth,\tl@runningwidth);
			\useasboundingbox (0,-1.5ex)
				rectangle (\hintscolumnwidth,1ex);
			\fill [#1] (0,0)
				++(\tl@startfraction*\hintscolumnwidth,0pt)
				rectangle (\tl@endfraction*\hintscolumnwidth,\tl@width);
			\node[#1,tl@singleyear,anchor=north] at (\tl@startfraction*\hintscolumnwidth/2+\tl@endfraction*\hintscolumnwidth/2,0pt) {#4};
		}
	}
	{#5}{#6}{#7}{#8}{#9}%
}


%-------------------------------------------
%   Define tldoublelabelcventry
%-------------------------------------------

\pgfkeys{
	/tldoublelabelcventry/.is family,
	/tldoublelabelcventry,
	default/.style={
		starta = 0,
		enda = 0,
		startb = 0,
		endb = 0,
		label = label,
	},
	starta/.estore in = \tlstarta,
	enda/.estore in = \tlenda,
	startb/.estore in = \tlstartb,
	endb/.estore in = \tlendb,
	label/.estore in = \tllabel,
}
\newcommand{\tldoublelabelcventry}[7][color1]{%
	\pgfkeys{/tldoublelabelcventry, default, #2}%
	\tl@formatstartyear{\tlstarta}
	\edef\tl@startfractiona{\tl@startfraction}
	\tl@formatendyear{\tlenda}
	\edef\tl@endfractiona{\tl@endfraction}
	\tl@formatstartyear{\tlstartb}
	\edef\tl@startfractionb{\tl@startfraction}
	\tl@formatendyear{\tlendb}
	\edef\tl@endfractionb{\tl@endfraction}
	\cventry{
		\tikz[baseline=0pt]{
			\fill [\tl@runningcolor] (0,0) rectangle (\hintscolumnwidth,\tl@runningwidth);
			\useasboundingbox (0,-1.5ex)
				rectangle (\hintscolumnwidth,1ex);
			\fill [#1] (0,0)
				++(\tl@startfractiona*\hintscolumnwidth,0pt)
				rectangle (\tl@endfractiona*\hintscolumnwidth,\tl@width);
			\fill [#1] (0,0)
				++(\tl@startfractionb*\hintscolumnwidth,0pt)
				rectangle (\tl@endfractionb*\hintscolumnwidth,\tl@width);
			\pgfmathsetmacro\pos{(\tl@startfractiona+\tl@startfractionb+\tl@endfractiona+\tl@endfractionb)/4*\hintscolumnwidth}
			\node[#1,tl@singleyear,anchor=north] at (\pos pt,0pt) {\tllabel};
		}
	}
	{#3}{#4}{#5}{#6}{#7}%
}

\makeatother

%-------------------------------------------
%   Theme related colors
%-------------------------------------------

% color0      main default color, normally left to black
% color1      primary scheme color
% color2      secondary scheme color

%----------------------------------------
%   General configuration
%----------------------------------------

% reduce underscores size
\renewcommand{\_}{\textscale{.5}{\textunderscore}}

%-------------------------------------------
%   Lengths modifications
%-------------------------------------------

\AtBeginDocument{%
    \setlength{\quotecustomwidth}{0.85\textwidth}
    %\setlength{\quotewidth}{0.65\textwidth}
    %\setlength{\separatorcolumnwidth}{0.025\textwidth}
    %\setlength{\maincolumnwidth}{\textwidth}
    %\settowidth{\listitemsymbolwidth}{\listitemsymbol}
    %\sethintscolumnlength{<length>} % <length> is the desired length in a unit LaTeX understands
    %\sethintscolumntowidth{<string>} % where <string> is a string of the desired length (usually, the longest string that has to appear in the column)
}

%-------------------------------------------
%   Theme font and style redefinition
%-------------------------------------------

\AtBeginDocument{%
    % fonts
    \renewcommand*{\titlefont}{\fontsize{32}{14}\mdseries\upshape}
    \renewcommand*{\namefont}{\fontsize{19}{14}\mdseries\upshape}
    %\renewcommand*{\addressfont}{\normalsize\mdseries\upshape}
    %\renewcommand*{\quotefont}{\large\slshape}
    %\renewcommand*{\sectionfont}{\Large\bfseries\upshape}
    %\renewcommand*{\subsectionfont}{\large\upshape\fontseries{sb}\selectfont}
    %\renewcommand*{\hintfont}{\bfseries}

    % styles
    \renewcommand*{\titlestyle}[1]{{\titlefont\textcolor{color1}{#1}}}
    \renewcommand*{\namestyle}[1]{{\namefont\textcolor{color2!85}{#1}}}
    %\renewcommand*{\addressstyle}[1]{{\addressfont\textcolor{color2}{#1}}}
    %\renewcommand*{\quotestyle}[1]{{\quotefont\textcolor{color1}{#1}}}
    %\renewcommand*{\sectionstyle}[1]{{\sectionfont\textcolor{color1}{#1}}}
    %\renewcommand*{\subsectionstyle}[1]{{\subsectionfont\textcolor{color1}{#1}}}
    %\renewcommand*{\hintstyle}[1]{{\hintfont\textcolor{color0}{#1}}}
}

%-------------------------------------------
%   hyperref colors
%-------------------------------------------

\AtEndPreamble{%
    \hypersetup{
        colorlinks=true,
        %linkcolor=red,
        %anchorcolor=black,
        citecolor=green,
        %filecolor=cyan,
        %menucolor=red,
        urlcolor=color1,
        linktoc=all
    }
}

%-------------------------------------------
%   Symbols from FontAwesome
%-------------------------------------------

\AtBeginDocument{%
    \renewcommand*{\addresssymbol}       {\faIcon{location-dot}\ }
    \renewcommand*{\mobilephonesymbol}   {\faIcon{mobile-screen}\ }
    \renewcommand*{\fixedphonesymbol}    {NONE\ }
    \renewcommand*{\faxphonesymbol}      {NONE\ }
    \renewcommand*{\emailsymbol}         {\faIcon[regular]{envelope}\ }
    \renewcommand*{\homepagesymbol}      {\faIcon{globe}\ }
    \renewcommand*{\linkedinsocialsymbol}{\faIcon{linkedin}\ }
    \renewcommand*{\twittersocialsymbol} {\faIcon{twitter}\ }
    \renewcommand*{\githubsocialsymbol}  {\faIcon{github}\ }
}

%-------------------------------------------
%   Commands definitions
%-------------------------------------------

% Used to indicate the level in the languages section
\newcommand{\cvskill}[2]{%
\textcolor{color2}{\textbf{#1}}
\foreach \x in {1,...,5}{%
	\space{\ifnumgreater{\x}{#2}{\color{color1!30}}{\color{color1}}\faCircle}}%
}

\newcommand{\cvitemcom}[2]{%
	\cvitemwithcomment{}{\listitemsymbol #1}{#2}%
}

% Colored href
\newcommand{\Chref}[3][blue]{\href{#2}{\color{#1}{#3}}}%

% Space to delete between the sections
\newcommand{\deletedSpace}{0mm}

% Space to delete before the quote
\newcommand{\quoteDeletedSpace}{0mm}

% C++
\newcommand{\Cpp}{%
  \texorpdfstring{%
    C\nolinebreak[4]\hspace{-.05em}%
    \raisebox{.4ex}{\smaller\smaller\bfseries ++}%
  }{C++}%
}

%-------------------------------------------
%   Age computation
%-------------------------------------------

\newcounter{birthday}
\newcounter{today}
\setmydatenumber{birthday}{1996}{06}{19}
\setmydatenumber{today}{\the\year}{\the\month}{\the\day}
\FPsub\result{\thetoday}{\thebirthday}
\FPdiv\myage{\result}{365.2425}
\FPtrunc\myage{\myage}{0}

%-------------------------------------------
%   Bibliography / publications
%-------------------------------------------

\bibliography{publications}

%-------------------------------------------
%   ModernCV theme
%-------------------------------------------

\moderncvtheme[blue]{banking}
%\moderncvtheme[blue]{classic}
%\renewcommand{\familydefault}{\sfdefault}
\nopagenumbers{}


% Space to delete between the sections
\renewcommand{\deletedSpace}{0mm}

% Space to delete before the quote
\renewcommand{\quoteDeletedSpace}{0mm}

%----------------------------------------------------------------------------------------
%   PERSONAL DATA
%----------------------------------------------------------------------------------------

\preinfo{%
    \faUser\ \myage{} ans
}
\firstname{Maxime}
\familyname{Pinard}
\title{Développeur \Cpp{}}
\address{Croix, mobilité nationale}{}
\phone[mobile]{+33~(0)6~87~92~55~09}
\email{maxime.pin@live.fr}
%\homepage{maxime.pinard.info}
\social[orcid]{0000-0001-6258-4687}
\social[linkedin]{maxime-pinard}
\social[github]{pinam45}
%\extrainfo{%
%    extra info
%}
\photo[60pt]{MPinard.jpg}
\quote{Passionné d'informatique et d'optimisation, je me spécialise en \Cpp{} et son écosystème.\\
    Mon projet personnel favori: \href{https://github.com/pinam45/dynamic_bitset}{dynamic\_bitset}}

%----------------------------------------------------------------------------------------
%   DOCUMENT BODY
%----------------------------------------------------------------------------------------

\begin{document}
    \maketitle

    %-------------------------------------------
    %   Work experience
    %-------------------------------------------

    \vspace*{0mm}
    \section{Expériences Professionnelles}
        \tlcventry{2020/6}{0}
            {Ingénieur de développement \Cpp{}}
            {IT-Finance, 6 ans (poste actuel)}
            {Wasquehal}
            {}
            {}
            \vspace{-5pt}
            \begin{itemize}
                \item Développement, optimisation et maintenance de systèmes de scans de marchés en \Cpp{}
                \item Implémentation de mémoires partagées POSIX avec chargement et modification en temps réel par flux de données
                \item Refonte et maintenance des systèmes de build avec CMake, et gestion des dépendances avec Conan 2
                %\item Mise en place du packaging interne et vendoring des libs externes avec CMake et Conan 2
                \item Mise en place des systèmes de CI sous Jenkins pour les tests, le code coverage et le packaging Debian
                \item Automatisations en Python pour le monitoring, le suivit des déploiements, la veille technologique\ldots
            \end{itemize}
            \vspace{5pt}
        \tllabelcventry{2019/9}{2020/2}{09/2019 - 02/2020}
            {Stagiaire recherche et développement logiciel}
            {Institut de Recherche en Informatique, Mathématiques, Automatique et Signal, 6 mois\hspace{-25pt}}
            {Mulhouse}
            {}
            {}
            \vspace{-5pt}
            \begin{itemize}
                \item Développement de métaheuristiques hybrides mémétiques pour le Unicost Set Cevering Problem\\
                en \Cpp{} avec recherche poussé de performances: amélioration des meilleurs résultats connu à l'époque
                \item Évaluations des performances sur le cluster HPC du méso-centre de Strasbourg et soumission à 2 conférences \cite{Pinard2020}
            \end{itemize}
            \vspace{5pt}
        \tldoublelabelcventry{%
                starta=2018/7,
                enda=2018/8,
                startb=2019/7,
                endb=2019/8,
                label={07-08/2018, 07-08/2019},
            }
            {Auxiliaire ambulancier}
            {Ambulances Phoenix, 2\times{}2 mois}
            {Héricourt}
            {Prise en charge et transport de patients}
            {}
        \tllabelcventry{2017/8}{2018/1}{08/2017 - 01/2018}
            {Stagiaire développement logiciel}
            {Direction Générale de l'Armement Maîtrise de l'Information, 6 mois}
            {Bruz}
            {}
            {}
            \vspace{-5pt}
            \begin{itemize}
                \item Travail sur GenDbg, un débogueur multi langages, OS et architectures:\\
                    \phantom{=}Développement en C du module de désassemblage pour les architectures MIPS et des tests unitaires associés
                \item Travail sur YaCo, plugin IDA Pro de rétro-ingénieurie collaborative utilisant libgit2:\\
                    \phantom{=}Portage en \Cpp{}, amélioration de la gestion interne git et de la capture des évènements IDA
            \end{itemize}
            \vspace{5pt}
        \tldoublelabelcventry{%
                starta=2015/2,
                enda=2015/2,
                startb=2016/2,
                endb=2016/2,
                label={02/2015, 02/2016},
            }
            {Animateur}
            {JAB France, 2\times{}1 semaine}
            {Evolène / Contamines-Montjoie}
            {Camp de ski de 50 jeunes, équipe de 20 animateurs}
            {}
        \tllabelcventry{2015/1}{2015/2}{01/2015}
            {Stagiaire}
            {Souchier SAS, 4 semaines}
            {Héricourt}
            {Jointage et montage d'appareils de désenfumage}
            {}

    %-------------------------------------------
    %   Education and qualifications
    %-------------------------------------------

    \vspace*{\deletedSpace}
    \section{Cursus}
        \tlcventry{2018/9}{2020/2}
            {Maîtrise en informatique (en double diplôme avec l'UTBM)}
            {Université du Québec à Chicoutimi}
            {Saguenay, Canada}
            {\smaller\textit{Bac\raisebox{.15ex}{\smaller+}5}}
            {}
        \tlcventry{2016/9}{2020/2}
            {Diplôme d'ingénieur en informatique, spécialité imagerie}
            {Université de Technologie de Belfort-Montbéliard}
            {Belfort}
            {\smaller\textit{Bac\raisebox{.15ex}{\smaller+}5}}
            {}
        \tlcventry{2014/9}{2016/8}
            {Diplôme d'études universitaires de technologie}
            {Université de Technologie de Belfort-Montbéliard}
            {Sèvenans}
            {\smaller\textit{Bac\raisebox{.15ex}{\smaller+}2}}
            {}
        % \tlcventry{2011/9}{2014/8}
        %     {Baccalauréat S, option SVT, spécialité Mathématiques}
        %     {Lycée Louis Aragon}
        %     {Héricourt}
        %     {\textit{Mention Bien}}
        %     {}

    %-------------------------------------------
    %   Computer skills
    %-------------------------------------------

    \vspace*{\deletedSpace}
    \section{Compétences informatiques}
        \cvline{Langages}{\Cpp{}, C, Python, SQL, Bash, Java, \LaTeX{}/Ti\textit{k}Z, CUDA, HTML/CSS, Assembleur\ldots}
        \cvline{Outils}{git, CMake, Make, Conan 2, Docker, GNU Tools, Clang Tools, Sanitizers, perf, Valgrind, Doxygen, gcov/lcov, Jenkins\ldots}
        \cvline{Méthodes}{Agile, Test unitaires / fonctionnels / intégration, TDD, Design patterns, UML, BPMN}
        \cvline{Métaheuristiques}{Algorithmes génétiques et mémétiques, recherche tabou, recuit simulé, colonies de fourmis, PSO}
        \cvline{Bibliothèques}{fmt, boost, PEGTL, spdlog, Catch2, GTest, google/benchmark, nlohmann/json, simdjson, SFML, ImGui, CLI11, cxxopts, inja, GLFW, OpenGL, OpenCL, OpenMP, MPI, Capstone, Keystone, GLM, libmpg123, TagLib\ldots}
        \cvitem{Administration système}{GNU/Linux (Debian), hébergement de services Docker, systèmes de monitoring et de backup}
        %\cvline{IDE}{CLion, Pycharm, IntelliJ IDEA, Visual Studio, Sublime Text}

    %-------------------------------------------
    \newpage
    %-------------------------------------------

    %-------------------------------------------
    %   Projects
    %-------------------------------------------

    %\vspace*{\deletedSpace}
    \section{Projets publics}
        \subsection{Personnels}
            \cvitemcom{Implémentation header-only d'un dynamic bitset et de ces test unitaires [\Cpp{}, Catch2]}
                {\href{https://github.com/pinam45/dynamic_bitset}{dynamic\_bitset}}
            \cvitemcom{Benchmarks de différentes implémentations de dynamic bitset (WIP) [\Cpp{}, google/benchmark]}
                {\href{https://github.com/pinam45/dynamic_bitset_benchmarks}{dynamic\_bitset\_benchmarks}}
            \cvitemcom{Lecteur de musique [\Cpp{}, ImGui, SFML, spdlog, libmpg123]}
                {\href{https://github.com/pinam45/MagicPlayer}{MagicPlayer}}
            \cvitemcom{Fonctions d'aide à la configuration des projets C/\Cpp{} (clang-format, ccache, sanitizers,\ldots) [CMake]}
                {\href{https://github.com/pinam45/CMutils}{CMutils}}
            \cvitemcom{Convertisseur de base, exemple ImGui/SFML [\Cpp{}, ImGui, SFML]}
                {\href{https://github.com/pinam45/BaseConverter}{BaseConverter}}
            \cvitemcom{Librairie graphique en console pour systèmes basé Unix et Windows [C]}
                {\href{https://github.com/pinam45/ConsoleControl}{ConsoleControl}}
            \cvitemcom{Dungeon crawler avec niveaux à génération procédurale [Java, JavaFx]}
                {\href{https://github.com/TiWinDeTea/Raoul-the-Game}{Raoul-the-Game}}
            \cvitemcom{Implémentation \LaTeX{}/Ti\textit{k}Z de la 1\textsuperscript{ère} et 4\textsuperscript{ème} de couverture des rapports de stage UTBM}
                {\href{https://github.com/pinam45/utbm-latex-internship-report-covers}{utbm-latex-internship-report-covers}}
            \cvitemcom{Implémentation \LaTeX{}/Ti\textit{k}Z Beamer du theme pour présentations UTBM}
                {\href{https://github.com/pinam45/utbm-beamer-theme}{utbm-beamer-theme}}
            % \cvitemcom{Jeu type Snake multijoueur (réseau local) [\Cpp{}, SFML]}
            %     {\href{https://github.com/TiWinDeTea/PapraGame}{PapraGame}}
        \subsection{Recherche}
            \cvitemcom{Calculateur d'hyperplans dans des géométries finies de dimension 4 (orienté performances) [\Cpp{}] \cite{Boulmier2019}}
                {\href{https://github.com/Lomadriel/HyperplaneFinder}{HyperplaneFinder}}
        \subsection{Stages}
            \cvitemcom{Programme de résolution du Set Cover Problem [\Cpp{}, dynamic\_bitset, spglog, json, cxxopts] \cite{Pinard2020}}
                {\href{https://github.com/pinam45/USCP}{USCP}, \href{https://github.com/pinam45/UTBM_ST50_Rapport_de_stage_IRIMAS/releases/latest}{Rapport de stage}}
            \cvitemcom{Programme de traitement des résultats du Set Cover Problem [\Cpp{}, spglog, json, cxxopts, inja] \cite{Pinard2020}}
                {\href{https://github.com/pinam45/USCP}{USCP}, \href{https://github.com/pinam45/UTBM_ST50_Rapport_de_stage_IRIMAS/releases/latest}{Rapport de stage}}
            \cvitemcom{Plugin Hex-Rays IDA de reverse-engineering collaboratif [\Cpp{}, GTest, IDA]}
                {\href{https://github.com/DGA-MI-SSI/YaCo}{YaCo}, \href{https://github.com/pinam45/UTBM_ST40_Rapport_de_stage_DGA/releases/latest}{Rapport de stage}}
            \cvitemcom{Module d'assemblage/désassemblage pour code assembleur MIPS [C, Capstone, Keystone, GTest]}
                {\href{https://github.com/pinam45/UTBM_ST40_Rapport_de_stage_DGA/releases/latest}{Rapport de stage}}
        \subsection{Universités}
            \cvitemcom{Algorithmes pour la résolution du Set Cover Problem [\Cpp{}, spdlog, dynamic\_bitset]}
                {\href{https://github.com/pinam45/UQAC_8INF870_Set_cover_problem}{Set\_cover\_problem}}
            \cvitem{}{~~~> Recherche exhaustive, Greedy, Branch and Bound, Recuit Simulé, Algorithme Génétique et Mémétique}
            \cvitemcom{Construction et rendu d'un objet 3D paramétrique avec OpenGL [\Cpp{}, GLFW, ImGui]}
                {\href{https://github.com/pinam45/UTBM\_IN55\_ParametricObjectsConstruction}{ParametricObjectsConstruction}}
            \cvitemcom{Clone de MiniMetro: simulateur de gestion de métros [Java, JavaFx]}
                {\href{https://github.com/TiWinDeTea/MagicMetro}{MagicMetro}}
            \cvitemcom{Jeux de Pogo avec IA MinMax/AlphaBeta [\Cpp{}, ConsoleControl]}
                {\href{https://github.com/pinam45/UTBM_IA41_Pogo}{Pogo}}
            %\cvlistitem{Serveur de stockage de fichiers multi-utilisateurs [\Cpp{}, SFML]}

    %-------------------------------------------
    %   Publications
    %-------------------------------------------

    \vspace*{\deletedSpace}
    \printbibliography[title=Publications]{}
    \vspace*{-2mm}

    %-------------------------------------------
    %   Languages
    %-------------------------------------------

    \vspace*{\deletedSpace}
    \section{Langues}
        \cvdoubleitem{\cvskill{5} \bfseries Français}{langue maternelle}{\cvskill{1} \bfseries Espagnol}{vagues souvenirs}
        \cvdoubleitem{\cvskill{4} \bfseries Anglais}{usage courant et professionnel}{\cvskill{0} \bfseries Japonais}{j'ai tout oublié}

    %-------------------------------------------
    %   Others qualifications
    %-------------------------------------------

    \vspace*{\deletedSpace}
    \section{Certifications autres}
        \cvlistitem{Permis de conduire}{}
        \cvlistitem{Prévention et Secours Civiques de niveau 1 (PSC1)}{}

    %-------------------------------------------
    %   Interests
    %-------------------------------------------

    \vspace*{\deletedSpace}
    \section{Centres d'intérêts}
        \cvline{L'informatique}
            {L'architecture des processeurs, la génération procédurale, la recherche opérationnelle, l'optimisation, la cryptographie, la sécurité et l'évolution de l'informatique quantique (veille technologique sur \href{http://links.pinard.io}{links.pinard.io})}
        \cvline{La cinéma}
            {200 à 250 films vu par an (diary sur \href{https://letterboxd.com/pinam45/diary/}{letterboxd})}
        \cvline{Les sciences}
            {Les mathématiques, la simulation de système physique, la physique quantique\ldots}
        \cvline{Le sport}
            {La randonnée, le ski et les sports d'hivers, la natation}
\end{document}
